\section{问题与路线}\label{sec:intro}

\subsection{任务与边界}

课程要求在一个颜色空间内对亚洲人脸进行分割，并对多幅人脸图像验证。我们把它落实为像素级的人脸区域分割：输入一张 $512\times512$ 的彩色图像，输出与原图同尺寸的二值掩膜，1 表示脸部，0 表示其他。评价的正类包括皮肤、鼻子、双眼、双眉、口腔和上下唇；眼镜作为忽略区域，不参加任何混淆计数。

有两个边界需要先说明。第一，肤色掩膜不等于人脸检测，更不等于身份识别。掩膜可能包含脖子、耳朵等皮肤区域，也不保证画面中只有一张脸；本报告不评价识别性能。第二，数据中的“亚洲人脸”来自实验者的外观筛选，没有官方人群属性标签；报告中的结论只对应这批图像及其固定划分，不外推到“所有亚洲人脸”。

\subsection{三个问题}

全文围绕三个问题展开。第一，这批人脸的肤色在 Cr/Cb 上是否足够集中，均值与标准差能否概括？第二，只靠颜色加两个空间规则，F1 能做到多少，提升来自哪一步？第三，它在什么图像条件下还成立，失败是参数问题还是模型问题？

对应地，我们经历了两个实验阶段的递进探索。实验三围绕纯颜色与基础连通域规则做了三次判断：第一次是颜色范围（$k$ 取多少），第二次是空间结构（最大连通域与填孔值多少），第三次是无约束补全强度（在填孔前加入 $41\times41$ 膨胀、填孔后再腐蚀能否补回内部缺口）。实验四则针对实验三暴露的“过早取最大连通域误删脸部碎片”与“头发脖颈粘连无法靠无约束形态学分离”两类核心瓶颈，引入基于灰度黑帽变换与双窗先验的双眼位置识别，由双眼几何关系生成正中与侧向内外嵌套双椭圆，结合低支持度样本放宽肤色阈值（$k=2.00$）重选及分区形态学连接，突破了纯颜色模型的等值线限制。

\subsection{证据状态与章节安排}

初版方案在固定配置后只运行一次测试，得到本报告唯一的独立测试结果。实验三的 $41\times41$ 补全尝试来自测试图 00570 的误差，因此在同一批测试图上的数字按开发后再评价报告；实验四的双眼定位与双椭圆约束形态学优化则严格固定在 32 张验证集（开发集）上进行多轮消融与验证，各阶段结果在表与图中均明确标注所属集合与证据状态，不相互混用。

第~\ref{sec:data}~章说明数据、标签与评价口径；第~\ref{sec:model}~章给出颜色模型与实验三初版方案；第~\ref{sec:test}~章报告独立测试并归纳三类失败；第~\ref{sec:completion}~章接第三类失败，做无约束形态学补全尝试并再次评价；第~\ref{sec:boundary}~章讨论纯颜色模型的能力边界与适用范围；第~\ref{sec:lab4}~章系统报告实验四的双眼定位、双椭圆空间约束、侧向椭圆候选、低支持放宽阈值重建与分区形态学连接优化；第~\ref{sec:conclusion}~章总结全文。

在独立测试集上，实验三初版 F1 为 70.69\%，$41\times41$ 补全后为 70.70\%，说明无空间锚点的形态学只能改变召回与精确率的分配；而在验证集上，实验四将全局 F1 从实验三的 72.70\% 大幅提升至 82.89\%（IoU 从 57.11\% 提升至 70.79\%），证明引入五官位置与双椭圆几何先验能够实质性跨越纯颜色分割的性能瓶颈。
